\documentclass[10pt,a4paper]{amsart}
\usepackage[margin=19mm]{geometry}
\usepackage{amsmath,amssymb,mathtools}
\usepackage[scr=rsfs,cal=euler]{mathalfa}
\usepackage{xfrac}
\usepackage[unicode,hidelinks,bookmarks=false]{hyperref}
\newcommand{\ie}{i.e.\ }
\newcommand{\cf}{cf.\ }
\newcommand{\resp}{respectively\ }
\newcommand{\Subsection}{\subsection}
\providecommand{\nobreakdash}{\nobreak\hbox{-}}
\setlength{\emergencystretch}{2em}
\multlinegap0pt
\usepackage{bm}
\usepackage{bbm}
\usepackage{dsfont}
\usepackage{xspace}
\usepackage{cancel}
\usepackage{mathtools}
\usepackage{amsthm}
\makeatletter
\@ifundefined{remark}{\newtheorem{remark}{Remark}}{}
\makeatother
\title[Extensible membranes in inviscid flow: aerodynamics and sing]{Extensible membranes in inviscid flow: aerodynamics and singular limits}
\author{Yu Jun Loo, Silas Alben}
\date{Source: arXiv:2609.16484v1 (CC BY 4.0)}
\begin{document}
\maketitle

\noindent\emph{Source and licence.} Extensible membranes in inviscid flow: aerodynamics and singular limits. Version arXiv:2609.16484v1 (CC BY 4.0), distributed under the Creative Commons Attribution 4.0 International licence, as recorded in the arXiv licence metadata for this exact version and shown on the abstract page https://arxiv.org/abs/2609.16484v1. Full rights notice: licenses/arxiv-2609.16484-CC-BY-4.0.txt.\par\smallskip

\noindent\textit{Editorial scope.} Counted unit: the Remark whose source label is \texttt{None} in the article (source0.tex lines 1321--1323, source label \texttt{None}), delivered together with the article's own complete original proof of it (source0.tex lines 1324--1869, one \texttt{proof} environment of 546 lines). No mathematics is written, repaired or re-derived: statement and proof are the article's own text line for line. Required context retained uncounted: positive-R2-evolution (lines 1236--1243). Every definition, display and label of those lines is retained unchanged. Omitted from this package and not redistributed: 1--1235, 1244--1320, 1870--2363. Nothing inside a retained excerpt is omitted or altered: every retained body is delivered exactly as the article's lines are written, so no omission receipt of any kind accompanies this package. Citations: the retained excerpts carry no citation command at all, so no bibliography entry of the article is transcribed and no reference is redistributed. Reference adaptation: none was needed and none was made; every reference inside the delivered unit resolves inside it, so no unresolved reference or citation is produced. Printed slips kept literal: no printed word, symbol, hypothesis or number of the counted statement or of its proof was altered. Layout and class: the article's own class is \texttt{article} with packages including authblk, tabto, relsize, graphicx, tikz, float, amsmath, amssymb, amsthm, amsfonts, adjustbox, algpseudocode, algorithm, comment; the wrapper is the lane's pinned \texttt{amsart} editorial wrapper at 10pt on A4, which restates the theorem-like environments the retained text needs and adds the article's own macro definitions that the retained text needs (none), with extra wrapper packages (bm, bbm, dsfont, xspace, cancel, mathtools). The delivered numbering is the unit's own. Assets: no retained span contains a figure environment, a graphics-inclusion command, a PDF-page inclusion, a table or a TikZ picture; no image file is copied into this package, the package declares no asset dependency and no image file is redistributed with it. Licence: version arXiv:2609.16484v1 of this article is distributed under the Creative Commons Attribution 4.0 International licence, as recorded in the arXiv licence metadata for this exact version and shown on the abstract page https://arxiv.org/abs/2609.16484v1. A full rights notice is delivered at licenses/arxiv-2609.16484-CC-BY-4.0.txt. Characters: every retained excerpt is pure ASCII, so the delivered engine, XeLaTeX with its default Latin Modern text font, reports no missing character.
\begin{equation}
R_1\partial_{tt}\zeta+R_2\partial_\alpha^4\zeta
=
R_3\partial_\alpha G(\partial_\alpha\zeta)+f,
\qquad
G(q)=\left(1-\frac{1}{|q|}\right)q.
\label{positive-R2-evolution}
\end{equation}

\begin{remark}
    In plain language, the condition that the membrane remain uniformly immersed is equivalent to asking that everywhere the membrane strain is bounded below by a constant $> -100\%$, where $-100\%$ corresponds to compression of a line segment to a point. 
\end{remark}

\begin{proof}
We first show that the nonlinear stretching force is locally Lipschitz near an immersed curve. Writing

$$
    G_i(q)=\left(1-\frac{1}{|q|}\right)q_i,
$$

direct differentiation gives

$$
    \frac{\partial G_i}{\partial q_j}(q)
    =
    \left(1-\frac{1}{|q|}\right)\delta_{ij}
    +\frac{q_iq_j}{|q|^3},
$$

and

$$
    \frac{\partial^2G_i}{\partial q_j\partial q_k}(q)
    =
    \frac{
        \delta_{ij}q_k+\delta_{ik}q_j+\delta_{jk}q_i
    }{|q|^3}
    -
    \frac{3q_iq_jq_k}{|q|^5}.
$$

Consequently, for every \(m>0\),

$$
    \left|
        \frac{\partial G_i}{\partial q_j}(q)
    \right|
    \leq C\left(1+\frac1m\right),
    \qquad
    \left|
        \frac{\partial^2G_i}{\partial q_j\partial q_k}(q)
    \right|
    \leq \frac{C}{m^2},
    \qquad |q|\geq m.
$$

Let

$$
    m_0=\min_{\alpha\in[-1,1]}
    |\partial_\alpha\zeta_0(\alpha)|>0.
$$

Since $H^2 \hookrightarrow C^1$, suppose that \(\zeta\) and \(\eta\) lie in an \(H^2\)-neighborhood of \(\zeta_0\) sufficiently small that

$$
    \|\partial_\alpha\zeta-\partial_\alpha\zeta_0\|_{L^\infty},
    \,
    \|\partial_\alpha\eta-\partial_\alpha\zeta_0\|_{L^\infty}
    \leq \frac{m_0}{2}.
$$

For \(0\leq\theta\leq1\), let

$$
    q_\theta
    :=
    \theta\partial_\alpha\zeta
    +(1-\theta)\partial_\alpha\eta
$$

lie on the line segment between $\partial_\alpha \zeta$ and $\partial_\alpha \eta$. Then

$$
\begin{aligned}
    |q_\theta|
    &\geq
    |\partial_\alpha\zeta_0|
    -
    \theta
    |\partial_\alpha\zeta-\partial_\alpha\zeta_0|
    -
    (1-\theta)
    |\partial_\alpha\eta-\partial_\alpha\zeta_0|\\
    &\geq \frac{m_0}{2}.
\end{aligned}
$$

Thus the entire line segment joining
\(\partial_\alpha\zeta\) and \(\partial_\alpha\eta\) remains in the region where the first two derivatives of \(G\) are bounded. By the fundamental theorem of calculus,

$$
\begin{aligned}
    &\frac{\partial G_i}{\partial q_j}
        (\partial_\alpha\zeta)
    -
    \frac{\partial G_i}{\partial q_j}
        (\partial_\alpha\eta)\\
    &\qquad=
    \int_0^1
    \frac{\partial^2G_i}{\partial q_j\partial q_k}
        (q_\theta)
    \left(
        \partial_\alpha\zeta_k
        -
        \partial_\alpha\eta_k
    \right)\,d\theta,
\end{aligned}
$$

and hence

$$
    \left|
    \frac{\partial G_i}{\partial q_j}
        (\partial_\alpha\zeta)
    -
    \frac{\partial G_i}{\partial q_j}
        (\partial_\alpha\eta)
    \right|
    \leq
    \frac{C}{m_0^2}
    |\partial_\alpha\zeta-\partial_\alpha\eta|.
$$

Using

$$
    \partial_\alpha G_i(\partial_\alpha\zeta)
    =
    \frac{\partial G_i}{\partial q_j}
        (\partial_\alpha\zeta)
    \partial_\alpha^2\zeta_j,
$$

we obtain

$$
\begin{aligned}
    &\partial_\alpha G_i(\partial_\alpha\zeta)
    -
    \partial_\alpha G_i(\partial_\alpha\eta)\\
    &\quad=
    \frac{\partial G_i}{\partial q_j}
        (\partial_\alpha\zeta)
    \left(
        \partial_\alpha^2\zeta_j
        -
        \partial_\alpha^2\eta_j
    \right)\\
    &\qquad+
    \left[
        \frac{\partial G_i}{\partial q_j}
            (\partial_\alpha\zeta)
        -
        \frac{\partial G_i}{\partial q_j}
            (\partial_\alpha\eta)
    \right]
    \partial_\alpha^2\eta_j.
\end{aligned}
$$

Therefore,

$$
\begin{aligned}
    &
    \left\|
        \partial_\alpha G(\partial_\alpha\zeta)
        -
        \partial_\alpha G(\partial_\alpha\eta)
    \right\|_{L^2}\\
    &\qquad\leq
    C\left(1+\frac1{m_0}\right)
    \|\partial_\alpha^2\zeta-\partial_\alpha^2\eta\|_{L^2}
    +
    \frac{C}{m_0^2}
    \|\partial_\alpha\zeta-\partial_\alpha\eta\|_{L^\infty}
    \|\partial_\alpha^2\eta\|_{L^2}.
\end{aligned}
$$

The one-dimensional Sobolev embedding \(H^1\hookrightarrow L^\infty\) now gives

\begin{equation}
    \left\|
        \partial_\alpha G(\partial_\alpha\zeta)
        -
        \partial_\alpha G(\partial_\alpha\eta)
    \right\|_{L^2}
    \leq
    C(m_0,M)\|\zeta-\eta\|_{H^2},
    \label{lipschitz estimate}
\end{equation}

provided

$$
    \|\zeta\|_{H^2},\|\eta\|_{H^2}\leq M.
$$

It follows that

$$
    \mathcal N(\zeta)
    :=
    \frac{R_3}{R_1}
    \partial_\alpha G(\partial_\alpha\zeta)
$$

is locally Lipschitz from \(H^2\) into \(L^2\) in a neighborhood of every immersed curve.
\newline

Let

$$
    A=\frac{R_2}{R_1},
$$

and recall that

$$
    \mathcal L=A\partial_\alpha^4
$$

denotes the positive self-adjoint beam operator with domain

$$
    D(\mathcal L)
    =
    \left\{
        v\in H^4((-1,1);\mathbb R^2):
        v(\pm1)=\partial_\alpha^2v(\pm1)=0
    \right\}.
$$

Its energy space is

$$
    V=D(\mathcal L^{1/2})
      =H^2((-1,1);\mathbb R^2)\cap H_0^1((-1,1);\mathbb R^2),
$$

and the norm \(\|\mathcal L^{1/2}v\|_{L^2}\) is equivalent to \(\|v\|_{H^2}\) on \(V\).

Using the functional calculus for the positive self-adjoint operator \(\mathcal L\), define

$$
    C(t)=\cos(t\mathcal L^{1/2}),
    \qquad
    S(t)=\mathcal L^{-1/2}\sin(t\mathcal L^{1/2}).
$$

These are the two fundamental solution operators for the linear beam equation

$$
    \partial_{tt}\zeta+\mathcal L\zeta=0.
$$

One may think of $\cos$ and $\sin$ here as an infinite polynomial series. More precisely, \(C(t)\zeta_0\) is the solution with initial displacement \(\zeta_0\) and zero initial velocity, while \(S(t)v_0\) is the solution with zero initial displacement and initial velocity \(v_0\). Thus the solution of the linear initial-value problem is

$$
    \zeta_{\mathrm{lin}}(t)
    =
    C(t)\zeta_0+S(t)v_0.
$$

It is easy to verify (for example with the Fourier transform) that the propagators satisfy

$$
    \|C(t)v\|_V\leq\|v\|_V,
    \qquad
    \|\partial_tC(t)v\|_{L^2}\leq\|v\|_V,
$$

for \(v\in V\), and

$$
    \|S(t)w\|_V\leq\|w\|_{L^2},
    \qquad
    \|\partial_tS(t)w\|_{L^2}\leq\|w\|_{L^2},
$$

for \(w\in L^2\).
\newline

With

$$
    \mathcal N(\zeta)
    =
    \frac{R_3}{R_1}
    \partial_\alpha G(\partial_\alpha\zeta),
$$

equation \eqref{positive-R2-evolution} becomes

$$
    \partial_{tt}\zeta+\mathcal L\zeta
    =
    \mathcal N(\zeta)+\frac{1}{R_1}f.
$$

By Duhamel's principle, $\zeta$ satisfies the implicit relation
\begin{equation}
\begin{aligned}
\zeta(t)
&=
C(t)\zeta_0+S(t)v_0+
\int_0^t
S(t-\tau)
\left(
\mathcal N(\zeta(\tau))
+\frac{1}{R_1}f(\tau)
\right)\,d\tau\ \\
&=
\zeta_{\mathrm{lin}}(t)
+
\int_0^t
S(t-\tau)
\left(
\mathcal N(\zeta(\tau))
+\frac{1}{R_1}f(\tau)
\right)\,d\tau.
\end{aligned}
\label{positive-R2-duhamel}
\end{equation}

We now close the immersion bootstrap. By the Sobolev embedding theorem, there exists positive constant \(C_E\) such that

$$
    \|\partial_\alpha v\|_{L^\infty}
    \leq C_E\|v\|_{H^2},
    \qquad v\in V.
$$

Choose $\rho$ such that

$$
    0<\rho<\frac{m_0}{2C_E},
    \qquad
    m_0=
    \min_{\alpha\in[-1,1]}
    |\partial_\alpha\zeta_0(\alpha)|.
$$

Consider the closed ball

$$
    \mathcal B_{\tau_f,\rho}
    =
    \left\{
        \zeta\in C([0,\tau_f];H^2):
        \begin{array}{l}
        (i) \ \ \ \zeta\text{ satisfies the fixed endpoint conditions},\\[1mm]
        (ii)\ \ \zeta(0)=\zeta_0,\\[1mm]
        (iii)\ \displaystyle
        \sup_{0\leq t\leq \tau_f}
        \|\zeta(t)-\zeta_0\|_{H^2}\leq\rho
        \end{array}
    \right\}.
$$

For every \(\zeta\in\mathcal B_{\tau_f,\rho}\),

$$
\begin{aligned}
    |\partial_\alpha\zeta(\alpha,t)|
    &\geq
    |\partial_\alpha\zeta_0(\alpha)|
    -
    |\partial_\alpha(\zeta-\zeta_0)(\alpha,t)|\\
    &\geq
    m_0-C_E\|\zeta(t)-\zeta_0\|_{H^2}\\
    &\geq\frac{m_0}{2}.
\end{aligned}
$$

Every curve in \(\mathcal B_{\tau_f,\rho}\) is therefore uniformly immersed, and the local Lipschitz estimate
\eqref{lipschitz estimate} holds uniformly on this set.
\newline

Define the fixed-point map

$$
\begin{aligned}
    (\Phi\zeta)(t)
    &=
    \zeta_{\mathrm{lin}}(t)\\
    &\quad+
    \int_0^t
    S(t-\tau)
    \left(
        \mathcal N(\zeta(\tau))
        +\frac{1}{R_1}f(\tau)
    \right)\,d\tau,
\end{aligned}
$$

and observe that 
\begin{equation}
\begin{aligned}
    \sup_{0\leq t\leq \tau_f}
    \|\Phi\zeta(t)-\zeta_0\|_{H^2}
    \leq{}&
    \sup_{0\leq t\leq \tau_f}
    \|\zeta_{\mathrm{lin}}(t)-\zeta_0\|_{H^2}\\
    &+
    C\int_0^{\tau_f}
    \left(
        \|\mathcal N(\zeta(\tau))\|_{L^2}
        +\frac{1}{R_1}\|f(\tau)\|_{L^2}
    \right)\,d\tau.
    \label{fixed point estimate}
\end{aligned}
\end{equation}

Since

$$
    \zeta_{\mathrm{lin}}(t)\longrightarrow\zeta_0
    \quad\text{in }H^2
    \qquad\text{as }t\to0,
$$

and \(\mathcal N\) is bounded on \(\mathcal B_{\tau_f,\rho}\), by shrinking \(\tau_f\) we may ensure that the right-hand side of equation \eqref{fixed point estimate} is smaller than \(\rho\). Hence \(\Phi\) maps \(\mathcal B_{\tau_f,\rho}\) into itself.
\newline

If \(\zeta,\eta\in\mathcal B_{\tau_f,\rho}\), then

$$
\begin{aligned}
    \sup_{0\leq t\leq \tau_f}
    \|\Phi\zeta(t)-\Phi\eta(t)\|_{H^2}
    &\leq
    C\int_0^{\tau_f}
    \|
        \mathcal N(\zeta(\tau))
        -
        \mathcal N(\eta(\tau))
    \|_{L^2}\,d\tau\\
    &\leq
    \tau_fC(m_0,\rho,\|\zeta_0\|_{H^2})
    \sup_{0\leq t\leq \tau_f}
    \|\zeta(t)-\eta(t)\|_{H^2}.
\end{aligned}
$$

By shrinking \(\tau_f\) again, we may ensure that \(\Phi\) is a contraction. Banach's fixed-point theorem therefore gives a unique solution of \eqref{positive-R2-duhamel}. Because the fixed point lies in \(\mathcal B_{\tau_f,\rho}\),

$$
    |\partial_\alpha\zeta(\alpha,t)|
    \geq\frac{m_0}{2}
$$

throughout its interval of existence. Differentiating the Duhamel formula in time and using

$$
    \partial_tS(t)=C(t)
$$

also gives

$$
    \partial_t\zeta\in C([0,\tau_f];L^2).
$$

Thus we obtain local existence and uniqueness of $\zeta$ on the interval $[0,\tau_f]$. where $\tau_f$ depends on $m_0$, $\|\zeta_0\|_{H^2}$ and $\|\zeta_1\|_{L^2}$ and $\|f\|_{L^2}$ only.
\newline 

It remains to show that this solution may be continued on any interval where $\zeta$ remains an immersion. To see this, suppose that the maximal existence time \(T_{\max}\) is finite but that, for some \(m>0\),

$$
    |\partial_\alpha\zeta(\alpha,t)|\geq m
$$

for every \(\alpha\in[-1,1]\) and \(0\leq t<T_{\max}\). Since

$$
    \partial_\alpha G_i(\partial_\alpha\zeta)
    =
    \frac{\partial G_i}{\partial q_j}
        (\partial_\alpha\zeta)
    \partial_\alpha^2\zeta_j
$$

and

$$
    \left|
        \frac{\partial G_i}{\partial q_j}(q)
    \right|
    \leq C_m
    \qquad\text{whenever }|q|\geq m,
$$

we have

$$
    \|\mathcal N(\zeta)\|_{L^2}
    \leq C_m\|\zeta\|_{H^2}.
$$

The Duhamel estimates therefore imply

$$
\begin{aligned}
    \|\zeta(t)\|_{H^2}
    +\|\partial_t\zeta(t)\|_{L^2}
    \leq{}&
    C\left(
        \|\zeta_0\|_{H^2}
        +\|v_0\|_{L^2}
    \right)\\
    &+
    C_m\int_0^t
    \|\zeta(\tau)\|_{H^2}\,d\tau\\
    &+
    \frac{C}{R_1}
    \int_0^t
    \|f(\tau)\|_{L^2}\,d\tau.
\end{aligned}
$$

Then by invoking Gronwall's inequality, we have 

$$
    \sup_{0\leq t<T_{\max}}
    \left(
        \|\zeta(t)\|_{H^2}
        +\|\partial_t\zeta(t)\|_{L^2}
    \right)<\infty.
$$

Hence by the same argument above, the solution may be continued for some additional time $\tau_f^* > 0$. which depends only on $m$, $\|\zeta(T_{max})\|_{H^2}$,$\|\partial_t\zeta|_{T_{max}}\|_{H^2}$ and $\|f\|_{L^2}$. This contradicts the maximality of \(T_{\max}\). Thus, if \(T_{\max}<\infty\), then

$$
    \liminf_{t\uparrow T_{\max}}
    \min_{\alpha\in[-1,1]}
    |\partial_\alpha\zeta(\alpha,t)|=0.
$$

Hence, the solution persists for as long as the membrane remains uniformly immersed. Note that no assumption was made on the sign of

$$
    T=R_3\bigl(|\partial_\alpha\zeta|-1\bigr)
$$
in the argument.
\end{proof}

\end{document}
